\section{Discussion}
\label{sec:discussion}

\paragraph{Where the small model stands.} On the commands a car assistant mostly receives, the 1.5B student is statistically
indistinguishable from the 72B teacher while being 65 times cheaper in prompt length, about eight times faster on the same GPU
and small enough for a laptop CPU. Its consistent weakness is knowing what it \emph{cannot} do: unsupported features that
resemble supported ones (``honk'', ``lane-departure warning'', ``left turn signal'' $\rightarrow$ hazard lights at $p=0.99$)
are mapped to the nearest real function, often with high confidence. Data helps (contrastive refusals, more samples) and a threshold helps against uncertain mistakes, but
neither removes confidently wrong decisions; this matches the over-triggering reported for functional tokens by
\citet{jafarzadeh2026fixedkeys}. A hybrid design that forwards uncertain requests to a larger model when a connection is
available is the natural next step for a product.

\paragraph{Form sensitivity is a distillation artifact.} The student reacts to surface form that carries no meaning:
lower-casing only the first letter of the 230 clean-test sentences raised false refusals of the v5 model from 3 to 21, and a
lower-case
\vi{hỏng xe giữa đường rồi}{my car broke down} was refused while the capitalized version was understood. The teacher writes
clean, capitalized, formal sentences, and sequence-level distillation transfers only what appears in the data. The same cause
explains the gap between validation (97--99\%) and the independently generated tests (75--88\%), the lexical lock-in of R1, and
the need for short and noisy data in v6. Diversity has to be designed into generation (assigned personas, contexts, lengths
and vocabulary), because asking the teacher to ``write many varied sentences'' does not produce it.

\paragraph{Shortcuts appear wherever a feature is one-sided.} Three times a feature of sentence form ended up on one side of
the label distribution---imperative mood on calls (R1), stories on refusals (R4), long sentences on refusals rather than
clarifications (before v5)---and each time the model learned it as a rule. Adding hard examples in contrastive pairs across
labels was the single most effective data intervention we found. A test set of mostly short commands did not reveal the R4
shortcut; only a development set containing story-like \emph{doable} requests did.

\paragraph{Measurement discipline changed conclusions.} Several results looked like improvements and were not: R2's gain on
the ambiguous group was one lucky seed; the ``student beats 72B on ambiguity'' impression came from a 39-item group and
reversed on 130 items; the threshold's cost on clear commands was 6--10 times larger on test than on validation; 4-bit
quantization passed on validation and lost five points on implicit requests at test time. Freezing test sets before scoring,
training two seeds, testing paired outputs and writing criteria down in advance was cheap compared with the cost of drawing
these wrong conclusions.

\paragraph{When RL helps.} GRPO needs prompts on which the policy sometimes succeeds and sometimes fails, and benefits from a
reasoning step that offers several paths to sample. A functional-token model trained by distillation offers neither: its
decision is one token, and most training prompts are already solved. RL on fresh, harder prompts or a preference method on
explicitly paired outputs \citep{ross2025when2call} would be better matched to this setting; we did not run either.

\section{Conclusion}
\label{sec:conclusion}

Nexa Drive shows that a 1.5B model with functional tokens, distilled from an open 72B teacher with about US\$100 of rented GPU
time, can serve as a Vietnamese in-car assistant on a laptop CPU: it handles clear and implicit requests about as well as the
prompted teacher, asks instead of guessing when a request is only a complaint, and is wrapped in code that keeps wrong
decisions away from the car. The project's main lessons are about data and measurement rather than architecture: hard examples
must be added in contrastive pairs, synthetic validation data hides the weaknesses that matter, and small test sets need
several seeds and paired tests before a difference is believed. The remaining weaknesses---refusing unsupported features,
sensitivity to quantization on ambiguous input, background speech, and false wake-ups---are measured and documented, and
define the next steps: real-user data, a hybrid on-device/cloud design and multi-microphone input.
